\section{Conclusion}

A finite-volume solver was developed to investigate the steady-state temperature distribution in a two-dimensional channel flow, accounting for convection, thermal diffusion, and viscous dissipation. The governing energy equation was discretized in space using central-difference approximations, resulting in a system of ordinary differential equations. Several time-integration methods were implemented to advance this system toward steady state, including the approximate factorization method (AFI), a sparse direct implicit method, and the explicit second- and fourth-order Runge--Kutta methods.

The numerical stability analysis showed that the explicit Runge--Kutta methods require sufficiently small time steps to maintain stability. The eigenvalue analysis provided estimates of their stability limits, which were then used to select suitable time steps for the solver comparison. The implicit methods allowed larger time steps, although their computational cost depended on the solution procedure and implementation.

The performance comparison showed that all four methods reached the prescribed residual tolerance of $10^{-12}$ on the $50\times20$ grid. Among the tested methods, the sparse direct solver achieved the shortest computation time, with a median runtime of approximately $0.14$ s. The AFI, RK2, and RK4 methods required approximately $6.28$, $8.01$, and $10.72$ s, respectively. Although RK4 required fewer iterations than RK2, its additional stages per iteration resulted in a higher total computation time. These results apply to the selected grid, time steps, convergence criterion, and implementations; they do not establish a general performance ranking for other configurations.

The temperature-field evolution illustrated the interaction between convection, thermal diffusion, and viscous dissipation during the approach to steady state. The parameter study further clarified their individual contributions. Increasing the Brinkman number strengthened viscous heating and produced larger departures from the initial linear temperature profile. Increasing the convection coefficient enhanced the downstream transport of the inlet temperature distribution, whereas reducing it allowed thermal diffusion and viscous dissipation to have a greater influence. In the Peclet-number study, the ratio $Br/Pe$ was kept constant so that the viscous-dissipation coefficient remained unchanged while the thermal-diffusion coefficient decreased with increasing Peclet number. The resulting temperature distributions demonstrated how the balance between these mechanisms determines the shape and magnitude of the temperature field.

Overall, the project provided a numerical framework for solving the two-dimensional energy equation and examining the effects of its governing parameters. The results also demonstrated that solver performance depends on the time-integration method, time-step selection, and computational implementation. Further work could include a more detailed grid-convergence study and comparisons of computational cost and accuracy over multiple grid resolutions. 
